\begin{textsample}{POS Dim 1 – 2018 – Score 24.00 – 2018/t000688.txt}  \label{ex:f1_pos_066}
I ' d like to introduce you to a tiny microorganism that you ' ve probably never heard of : its name is Prochlorococcus, and it ' s really an amazing little being.
For one thing, its ancestors changed \textbf{the} earth in ways that made it possible for us to \textbf{evolve}, and hidden in its genetic code is a blueprint that may inspire ways to reduce our dependency on \textbf{fossil} fuel.
But \textbf{the} most amazing thing is that there are three billion billion billion of these tiny cells on \textbf{the} \textbf{planet}, and we didn ' t know they existed until 35 years ago.
So to tell you their story, I need to first take you way back, four billion years ago, when \textbf{the} earth might have looked something like this.
There was no life on \textbf{the} \textbf{planet}, there was no oxygen in \textbf{the} atmosphere.
So what happened to change that \textbf{planet} into \textbf{the} one we \textbf{enjoy} today, teeming with life, teeming with plants and animals?
Well, in a word, photosynthesis.
About two and a half billion years ago, some of these ancient ancestors of Prochlorococcus \textbf{evolved} so that they could use solar energy and absorb it and split water into its component parts of oxygen and hydrogen.
And they used \textbf{the} chemical energy produced to draw CO2, \textbf{carbon} dioxide, out of \textbf{the} atmosphere and use it to build sugars and proteins and amino acids, all \textbf{the} things that life is made of.
And as they \textbf{evolved} and grew more and more over millions and millions of years, that oxygen accumulated in \textbf{the} atmosphere.
Until about 500 million years ago, there was enough in \textbf{the} atmosphere that larger organisms could \textbf{evolve}.
There was an \textbf{explosion} of life- forms, and, ultimately, we appeared on \textbf{the} scene.
While that was going on, some of those ancient photosynthesizers died and were compressed and buried, and became \textbf{fossil} fuel with sunlight buried in their \textbf{carbon} bonds.
They ' re basically buried sunlight in \textbf{the} form of \textbf{coal} and \textbf{oil}.
Today ' s photosynthesizers, their engines are descended from those ancient microbes, and they feed basically all of life on earth.
Your heart is beating using \textbf{the} solar energy that some plant processed for you, and \textbf{the} stuff your body is made out of is made out of CO2 that some plant processed for you.
Basically, we ' re all made out of sunlight and \textbf{carbon} dioxide.
Fundamentally, we ' re just hot air.
So as terrestrial beings, we ' re very familiar with \textbf{the} plants on land : \textbf{the} trees, \textbf{the} grasses, \textbf{the} pastures, \textbf{the} crops.
But \textbf{the} oceans are filled with billions of tons of animals.
Do you ever wonder what ' s feeding them?
Well there ' s an invisible pasture of microscopic photosynthesizers called phytoplankton that fill \textbf{the} upper 200 meters of \textbf{the} ocean, and they feed \textbf{the} entire open ocean ecosystem.
Some of \textbf{the} animals live among them and eat them, and others swim up to feed on them at night, while others sit in \textbf{the} deep and wait for them to die and settle down and then they chow down on them.
So these tiny phytoplankton, collectively, weigh less than one percent of all \textbf{the} plants on land, but annually they photosynthesize as much as all of \textbf{the} plants on land, including \textbf{the} Amazon rainforest that we consider \textbf{the} lungs of \textbf{the} \textbf{planet}.
Every year, they fix 50 billion tons of \textbf{carbon} in \textbf{the} form of \textbf{carbon} dioxide into their bodies that feeds \textbf{the} ocean ecosystem.
How does this tiny amount of biomass produce as much as all \textbf{the} plants on land?
Well, they don ' t have trunks and stems and flowers and fruits and all that to maintain.
All they have to do is grow and divide and grow and divide.
They ' re really lean little photosynthesis machines.
They really crank.
So there are thousands of different species of phytoplankton, come in all different shapes and sizes, all roughly less than \textbf{the} \textbf{width} of a human hair.
Here, I ' m showing you some of \textbf{the} more beautiful ones, \textbf{the} \textbf{textbook} versions.
I call them \textbf{the} charismatic species of phytoplankton.
And here is Prochlorococcus.
I know, it just looks like a bunch of schmutz on a microscope \textbf{slide}.
But they ' re in there, and I ' m going to reveal them to you in a minute.
But first I want to tell you how they were discovered.
About 38 years ago, we were playing around with a technology in my lab called flow cytometry that was developed for biomedical research for studying cells like cancer cells, but it turns out we were using it for this off-label purpose which was to study phytoplankton, and it was beautifully suited to do that.
And here ' s how it works : so you inject a sample in this tiny little capillary tube, and \textbf{the} cells go single file by a laser, and as they do, they scatter light according to their size and they emit light according to whatever pigments they might have, whether they ' re natural or whether you stain them.
And \textbf{the} chlorophyl of phytoplankton, which is green, emits red light when you shine \textbf{blue} light on it.
And so we used this instrument for several years to study our phytoplankton cultures, species like those charismatic ones that I showed you, just studying their basic cell \textbf{biology}.
But all that time, we thought, well wouldn ' t it be really cool if we could take an instrument like this out on a ship and just squirt seawater through it and see what all those diversity of phytoplankton would look like.
So I managed to get my hands on what we call a big rig in flow cytometry, a large, powerful laser with a money-back guarantee from \textbf{the} company that if it didn ' t work on a ship, they would take it back.
And so a young scientist that I was working with at \textbf{the} time, Rob Olson, was able to take this thing apart, put it on a ship, put it back together and take it off to \textbf{sea}.
And it worked like a charm.
We didn ' t think it would, because we thought \textbf{the} ship ' s vibrations would get in \textbf{the} way of \textbf{the} focusing of \textbf{the} laser, but it really worked like a charm.
And so we mapped \textbf{the} phytoplankton distributions across \textbf{the} ocean.
For \textbf{the} first time, you could look at them one cell at a time in real time and see what was going on — that was very exciting.
But one day, Rob noticed some faint signals coming out of \textbf{the} instrument that we dismissed as electronic noise for probably a year before we realized that it wasn ' t really behaving like noise.
It had some regular patterns to it.
To make a long story short, it was tiny, tiny little cells, less than one-one hundredth \textbf{the} \textbf{width} of a human hair that contain chlorophyl.
That was Prochlorococcus.
So remember this \textbf{slide} that I showed you?
If you shine \textbf{blue} light on that same sample, this is what you see : two tiny little red light-emitting cells.
Those are Prochlorococcus.
They are \textbf{the} smallest and most abundant photosynthetic cell on \textbf{the} \textbf{planet}.
At first, we didn ' t know what they were, so we called \textbf{the}"little greens."It was a very affectionate name for them.
Ultimately, we knew enough about them to give them \textbf{the} name Prochlorococcus, which means"primitive green berry."And it was about that time that I became so smitten by these little cells that I redirected my entire lab to study them and nothing else, and my loyalty to them has really paid off.
They ' ve given me a tremendous amount, including bringing me here.
So over \textbf{the} years, we and others, many others, have studied Prochlorococcus across \textbf{the} oceans and found that they ' re very abundant over wide, wide ranges in \textbf{the} open ocean ecosystem.
They ' re particularly abundant in what are called \textbf{the} open ocean gyres.
These are sometimes referred to as \textbf{the} deserts of \textbf{the} oceans, but they ' re not deserts at all.
Their deep \textbf{blue} water is teeming with a hundred million Prochlorococcus cells per liter.
If you crowd them together like we do in our cultures, you can see their beautiful green chlorophyl.
One of those test tubes has a billion Prochlorococcus in it, and as I told you earlier, there are three billion billion billion of them on \textbf{the} \textbf{planet}.
That ' s three octillion, if you care to convert.
And collectively, they weigh more than \textbf{the} human population and they photosynthesize as much as all of \textbf{the} crops on land.
They ' re incredibly important in \textbf{the} global ocean.
So over \textbf{the} years, as we were studying them and found how abundant they were, we thought, hmm, this is really strange.
How can a single species be so abundant across so many different \textbf{habitats}?
And as we isolated more into culture, we learned that they are different ecotypes.
There are some that are adapted to \textbf{the} high-light intensities in \textbf{the} \textbf{surface} water, and there are some that are adapted to \textbf{the} low light in \textbf{the} deep ocean.
In fact, those cells that live in \textbf{the} bottom of \textbf{the} sunlit zone are \textbf{the} most efficient photosynthesizers of any known cell.
And then we learned that there are some strains that grow optimally along \textbf{the} equator, where there are higher temperatures, and some that do better at \textbf{the} cooler temperatures as you go north and south.
So as we studied these more and more and kept finding more and more diversity, we thought, oh my God, how diverse are these things?
And about that time, it became possible to sequence their \textbf{genomes} and really look under \textbf{the} hood and look at their genetic makeup.
And we ' ve been able to sequence \textbf{the} \textbf{genomes} of cultures that we have, but also recently, using flow cytometry, we can isolate individual cells from \textbf{the} wild and sequence their individual \textbf{genomes}, and now we ' ve sequenced hundreds of Prochlorococcus.
And although each cell has roughly 2, 000 genes — that ' s one tenth \textbf{the} size of \textbf{the} human \textbf{genome} — as you sequence more and more, you find that they only have a thousand of those in common and \textbf{the} other thousand for each individual strain is drawn from an enormous gene pool, and it reflects \textbf{the} particular environment that \textbf{the} cell might have thrived in, not just high or low light or high or low temperature, but whether there are nutrients that limit them like nitrogen, phosphorus or \textbf{iron}.
It reflects \textbf{the} \textbf{habitat} that they come from.
Think of it this way.
If each cell is a \textbf{smartphone} and \textbf{the} apps are \textbf{the} genes, when you get your \textbf{smartphone}, it comes with these built-in apps.
Those are \textbf{the} ones that you can ' t delete if you ' re an iPhone person.
You press on them and they don ' t jiggle and they don ' t have x ' s.
Even if you don ' t want them, you can ' t get rid of them.
Those are like \textbf{the} core genes of Prochlorococcus.
They ' re \textbf{the} essence of \textbf{the} phone.
But you have a huge pool of apps to draw upon to make your phone custom- designed for your particular lifestyle and \textbf{habitat}.
If you travel a lot, you ' ll have a lot of travel apps, if you ' re into financial things, you might have a lot of financial apps, or if you ' re like me, you probably have a lot of weather apps, hoping one of them will tell you what you want to hear.
And I ' ve learned \textbf{the} last couple days in Vancouver that you don ' t need a weather app — you just need an umbrella.
So — So just as your \textbf{smartphone} tells us something about how you live your life, your lifestyle, reading \textbf{the} \textbf{genome} of a Prochlorococcus cell tells us what \textbf{the} pressures are in its environment.
It ' s like reading its diary, not only telling us how it got through its day or its week, but even its \textbf{evolutionary} history.
As we studied — I said we ' ve sequenced hundreds of these cells, and we can now project what is \textbf{the} total genetic size — gene pool — of \textbf{the} Prochlorococcus federation, as we call it.
It ' s like a superorganism.
And it turns out that projections are that \textbf{the} collective has 80, 000 genes.
That ' s four times \textbf{the} size of \textbf{the} human \textbf{genome}.
And it ' s that diversity of gene pools that makes it possible for them to dominate these large regions of \textbf{the} oceans and maintain their stability year in and year out.
So when I daydream about Prochlorococcus, which I probably do more than is healthy — I imagine them floating out there, doing their job, maintaining \textbf{the} \textbf{planet}, feeding \textbf{the} animals.
But also I inevitably end up thinking about what a masterpiece they are, finely tuned by millions of years of \textbf{evolution}.
With 2, 000 genes, they can do what all of our human ingenuity has not figured out how to do yet.
They can take solar energy, CO2 and turn it into chemical energy in \textbf{the} form of \textbf{organic} \textbf{carbon}, locking that sunlight in those \textbf{carbon} bonds.
If we could figure out exactly how they do this, it could inspire designs that could reduce our dependency on \textbf{fossil} fuels, which brings my story full circle. \textbf{The} \textbf{fossil} fuels that are buried that we ' re burning took millions of years for \textbf{the} earth to bury those, including those ancestors of Prochlorococcus, and we ' re burning that now in \textbf{the} blink of an eye on geological timescales. \textbf{Carbon} dioxide is increasing in \textbf{the} atmosphere.
It ' s a \textbf{greenhouse} gas. \textbf{The} oceans are starting to warm.
So \textbf{the} question is, what is that going to do for my Prochlorococcus?
And I ' m sure you ' re expecting me to say that my beloved microbes are doomed, but in fact they ' re not.
Projections are that their populations will expand as \textbf{the} ocean warms to 30 percent larger by \textbf{the} year 2100.
Does that make me happy?
Well, it makes me happy for Prochlorococcus of course — but not for \textbf{the} \textbf{planet}.
There are winners and losers in this global experiment that we ' ve \textbf{undertaken}, and it ' s projected that among \textbf{the} losers will be some of those larger phytoplankton, those charismatic ones which are expected to be reduced in numbers, and they ' re \textbf{the} ones that feed \textbf{the} zooplankton that feed \textbf{the} \textbf{fish} that we like to harvest.
So Prochlorococcus has been my muse for \textbf{the} past 35 years, but there are legions of other microbes out there maintaining our \textbf{planet} for us.
They ' re out there ready and waiting for us to find them so they can tell their stories, too.
Thank you.

% matched lemmas: biology, blue, carbon, coal, enjoy, evolution, evolutionary, evolve, explosion, fish, fossil, genome, greenhouse, habitat, iron, oil, organic, planet, sea, slide, smartphone, surface, textbook, the, undertake, width
\end{textsample}
